\section{Hyperbolic theories: LS, GL and GN}\label{sec:hyperbolic}
\subsection{A common dispersion equation}
For plane longitudinal motion the three classical generalized theories fit into one equation. Let $\tau_0$ denote the flux relaxation time (Lord--Shulman) or the thermal relaxation time (Green--Lindsay, in the notation of the later literature). Let $\tau^\dagger$ be the lag that multiplies the thermal-stress and dilatation coupling. Eliminating $\hat u$ and $\hat\theta$ gives
\begin{equation}
\bigl(\rho c_1^2K+\rho s^2\bigr)\Bigl(s\rho c\,(1+\tau_0s)+kK\Bigr)+\beta^2T_0\,s\,K\,(1+\tau^\dagger s)=0,\label{eq:unified}
\end{equation}
with $\tau^\dagger=0$ and $\tau_0=0$ for the classical theory, $\tau^\dagger=\tau_0$ for Lord--Shulman (this is \eqref{eq:disp} multiplied by $1+\tau_0s$), and $\tau^\dagger=\tau_1$, the second relaxation time, for Green--Lindsay. The unified form of Dhaliwal and Sherief \citep{Dhaliwal1980} makes the same point at the level of the field equations. The theories differ in which of the two coupling terms carries a time-derivative correction.

\subsection{Lord--Shulman theory}\label{sec:ls}
The Lord--Shulman (LS) theory joins Cattaneo's law \eqref{eq:cattaneo} to the energy balance \eqref{eq:en}. Eliminating the flux gives the temperature equation
\begin{equation}
k\nabla^2\theta=(1+\tau_0\partial_t)\bigl(\rho c\,\dot\theta+T_0\beta\,\dot\varepsilon_{kk}\bigr)-(1+\tau_0\partial_t)\rho r, \label{eq:LS}
\end{equation}
so the lagged derivative acts on \emph{both} the heat-capacity term and the dilatation-coupling term. The second of these is a modeling consequence of applying the Cattaneo law to the total heat flux, and it produces the difference from the Green--Lindsay equations.

\paragraph{Wave fronts.} Take $\omega\to\infty$ in \eqref{eq:unified}. The leading terms in $s$ give, with $c_T^2=k/(\rho c\tau_0)$ and $\varepsilon=\beta^2T_0/(\rho^2cc_1^2)$,
\begin{equation}
c_{\rm f}^4-c_{\rm f}^2\bigl[c_1^2(1+\varepsilon\,\tau^\dagger/\tau_0)+c_T^2\bigr]+c_1^2c_T^2=0, \label{eq:fronts}
\end{equation}
whose two roots $c_{\rm f}$ are the speeds of the two wave fronts, one elastic-type and one thermal-type, with $v_{\rm e}<v_{\rm t}$ when $c_T>c_1$ (the labels exchange for $c_T<c_1$). For LS, $\tau^\dagger=\tau_0$ and the equation is the standard high-frequency front equation of the theory; for GL it is \eqref{eq:fronts} with the factor $\tau_1/\tau_0$. The front equation \eqref{eq:fronts} was verified numerically against the roots of \eqref{eq:unified} at $\omega/\omega_*=2\times10^3$ for both theories (front-speed check). With $\varepsilon=0.1$ and $\tau_0=0.25/\omega_*$ ($c_T=2c_1$) the LS fronts are $0.9839\,c_1$ and $2.0327\,c_1$. With $\tau_1=0.5/\omega_*$ (admissible, since $\tau_1\ge\tau_0$, and giving $\tau_1/\tau_0=2$) the GL fronts are $0.9689\,c_1$ and $2.0643\,c_1$, in agreement with \eqref{eq:fronts} to four digits. Two features deserve recording.
\begin{enumerate}
\item The elastic front speed is \emph{not} the isothermal speed $c_1$, nor the adiabatic speed. Interaction with the thermal front displaces it, and the two speeds repel one another: their product is $c_1c_T$ and the sum of their squares is $c_1^2(1+\varepsilon\,\tau^\dagger/\tau_0)+c_T^2$, that is $c_1^2(1+\varepsilon)+c_T^2$ for LS. For LS, when $c_T=c_1$ the degeneracy is lifted by an amount proportional to $\sqrt\varepsilon$, $v^2=c_1^2[2+\varepsilon\pm\sqrt{\varepsilon(4+\varepsilon)}]/2\approx c_1^2(1\pm\sqrt\varepsilon)$.
\item The leading front speeds depend only on the high-frequency terms, so they are independent of $k$ except through $c_T$.
\end{enumerate}

The two fronts are sketched in \cref{fig:fronts}.

\begin{figure}[t!]\centering\small
\input{figs_tikz/fig3_fronts.tex}
\caption{Original schematic; the setting follows \citep{Lord1967}: two wave fronts in hyperbolic thermoelasticity -- elastic-type and thermal-type fronts repel in speed; the thermal front carries a damped tail.}\label{fig:fronts}
\end{figure}

\paragraph{Discontinuities and well-posedness.} The system is hyperbolic, so initial data with a jump produce solutions whose jumps travel along the fronts and decay exponentially along the propagation path. The heat-pulse problem of \cref{sec:transient} shows this. The system is well posed for $\tau_0>0$, and energy decays in bounded bodies under standard boundary conditions by the energy identity \eqref{eq:energy}. Existence of generalized solutions was established for one relaxation time by \citep{Bem1982} and for two by \citep{Bem1983}, and domain-of-influence results appear in \citep{Ignaczak1980}. The thermodynamic foundation via the Clausius--Duhem inequality is due to \citep{GurtinWilliams1966}.

\paragraph{Interpretation.} Two modeling questions matter: the physical size of $\tau_0$, and whether Cattaneo's law with a single exponential relaxation describes the frequency range of interest. Joseph and Preziosi \citep{Joseph1989,Joseph1990} review the evidence for and against. Large relaxation times have been reported for nonhomogeneous materials \citep{Kaminski1990,Mitra1995}, though the interpretation of such measurements is debated. For solids the strongest evidence of wave-like heat propagation remains second sound at cryogenic temperatures \citep{Ackerman1966,Jackson1970}. There the governing regime is the Guyer--Krumhansl hydrodynamic one \citep{GuyerKrumhansl1966II} rather than the simple Cattaneo law.

\subsection{Green--Lindsay theory}\label{sec:gl}
The Green--Lindsay (GL) theory \citep{Green1972} has the constitutive relations
\begin{equation}
\sigma_{ij}=C_{ijkl}\varepsilon_{kl}-\beta_{ij}\bigl(\theta+\tau_1\dot\theta\bigr),\qquad
k\theta_{,ii}=\rho c\bigl(\dot\theta+\tau_0\ddot\theta\bigr)+T_0\beta_{ij}\dot\varepsilon_{ij},\label{eq:GL}
\end{equation}
with $\tau_1\ge\tau_0>0$. In the notation of Green and Lindsay \citep{Green1972}, $\tau_1$ is the coefficient $\alpha$ of $\dot\theta$ in the stress relation, and $\tau_0=h/d$ is the ratio of the coefficients of $\ddot\theta$ and $\dot\theta$ in the energy equation. Their restrictions $d\alpha-h\ge0$, $h>0$ (Eqs.~(4.14) and~(4.16) of that paper) give the ordering stated. Heat conduction is Fourier's here, so the theory does not fit the flux template \eqref{eq:template}: the finite thermal speed comes from the inertia term $\tau_0\ddot\theta$ instead. Substitution in plane waves gives \eqref{eq:unified} with $\tau^\dagger=\tau_1$. Three differences from LS matter. First, the thermal equation is not of Cattaneo type, so its relation to kinetic arguments is less direct. Second, the entropy depends on the temperature rate, which puts a thermal-displacement-like quantity into the constitutive structure and makes care necessary in reading $\tau_0$ and $\tau_1$ as independent material constants. Third, solutions of the original theory can contain discontinuities at the wave fronts, and a modification with a strain-rate term was proposed to remove them \citep{Yu2018ModifiedGL}. Wave fronts follow \eqref{eq:fronts} with $\tau^\dagger=\tau_1$, and one consequence distinguishes GL from LS. At the wave fronts the coupling is stronger than in LS by the factor $\tau_1/\tau_0\ge1$. The temperature-rate term in the stress carries $\tau_1$, the thermal inertia carries $\tau_0$, and the strain-rate coupling in the energy equation carries no relaxation time.

\subsection{Green--Naghdi types I, II and III}\label{sec:gn}
The Green--Naghdi (GN) theories \citep{Green1991,GreenGN1992,Green1993} postulate an entropy balance with the thermal displacement $v$ ($\dot v=\theta$) as a primary field. In the template the three types are
\begin{equation}
\text{I: }\ q_i=-k\theta_{,i};\qquad
\text{II: }\ q_i=-k^*v_{,i};\qquad
\text{III: }\ q_i=-k\theta_{,i}-k^*v_{,i},
\end{equation}
with $G=k$, $k^*/s$ and $k+k^*/s$, respectively. Their uncoupled thermal equations are, with $\rho c=1$ for brevity,
\begin{equation}
\text{II: }\ \ddot\theta=c_*^2\nabla^2\theta,\quad c_*^2=k^*/\rho c;\qquad
\text{III: }\ \ddot\theta=c_*^2\nabla^2\theta+\alpha\nabla^2\dot\theta.
\end{equation}
\paragraph{Type II.} Thermal disturbances propagate as undamped waves of speed $c_*$, and with coupling the two fronts satisfy \eqref{eq:fronts} with $c_T\to c_*$ and $\tau^\dagger/\tau_0\to1$. The relation is \emph{exact at every frequency}, not merely asymptotically: the roots of the quadratic \eqref{eq:disp} with $G=k^*/s$ are non-dispersive, with $v^4-v^2[c_1^2(1+\varepsilon)+c_*^2]+c_1^2c_*^2=0$ at every $\omega$. This was confirmed numerically for $\varepsilon=0.1$ and $c_*=c_1$: the exact speeds $0.854309c_1$ and $1.170537c_1$ are the same at $\omega/\omega_*=0.5$, $2$ and $10$, and agree with the formula. The theory conserves energy, so there is no thermoelastic damping at all. It is the only member of the family with that property, and therefore the benchmark for energy-based checks.
\paragraph{Type III.} The type III equation is a damped wave equation with a Kelvin--Voigt-type damping term $\alpha\nabla^2\dot\theta$. Two consequences deserve emphasis. At \emph{low} frequency $G\simeq k^*/s$ and the thermal branch has the type-II speed $c_*$. At \emph{high} frequency $G\to k$ and the branch becomes diffusive, so its phase speed is unbounded (\Cref{fig:family}, magenta curve), exactly as in the classical theory. Type III is thus wave-like at low frequencies and diffusive at high frequencies, in contrast to LS. Its thermal signals have no finite propagation speed in the strict sense, since the strongly damped wave equation has a parabolic smoothing effect. In practice type III describes weakly damped thermal waves at moderate frequencies; it is not an ``improved Cattaneo law''. Green and Naghdi's statement that type III admits waves of finite speed refers to the intermediate-frequency speed $c_*=\sqrt{k^*/\rho c}$ of the undamped part. It does not contradict the high-frequency behavior above, where the phase speed grows without bound and there is no sharp front. The word ``speed'' therefore names two different quantities, and a statement about type III should say which is meant.
\paragraph{Uniqueness and stability.} Chandrasekharaiah proved uniqueness for type II \citep{Chandrasekharaiah1996} and gave complete solutions \citep{Chandrasekharaiah1997}. The theory without energy dissipation for anisotropic and micropolar media was developed in \citep{Iesan1998,Quintanilla1999,Ciarletta1999,Quintanilla2002}. Stability and decay results for type III belong to the mathematical theory of thermoelasticity \citep{Suh1998}. For the type III system, Zhang and Zuazua \citep{ZhangZuazua2003} report three results. For most domains the energy does not decay uniformly, under suitable conditions on the domain it decays exponentially, and in two space dimensions the energy of smooth solutions decays polynomially. Decay of type III is therefore not unconditional. Type II has drawn attention as a conservative reference model, and combined with nonlocal elasticity it has been analyzed for wave dispersion \citep{Sarkar2019}.

\subsection{What each attribution was checked against}\label{sec:provenance}

An attribution is a claim about a source, and it can be checked against that source or merely copied from it. \Cref{tab:provenance} records, for the attributions in this review that were checked at page level, what the check consisted of and, just as important, what it does not cover. Statements absent from the table were not checked at page level and are presented as reproduced, which is the weaker claim.

\begin{table}[H]\centering\small
\caption{Page-level provenance for attributions made in this review. `Covers' names the statement the check
supports; `does not cover' names the part of this review that the check leaves standing on the secondary
literature.}
\label{tab:provenance}
\begin{tabularx}{\textwidth}{@{}p{3.1cm}Y Y@{}}
\toprule
\textbf{Statement} & \textbf{Checked against} & \textbf{Covers / does not cover} \\
\midrule
Biot heat-conduction equation, \eqref{eq:en} & Biot (1956), printed p.~242, Eqs.~(3.3)--(3.4) &
  Covers the linear isotropic Fourier form with constant coefficients in a source-free body. Does not
  cover coupled boundary-value solutions, anisotropy, or temperature-dependent properties. \\
Green--Lindsay stress and entropy, \eqref{eq:GL} & Green and Lindsay (1972), primary text &
  Covers the temperature-rate dependent stress and entropy, the absence of a relaxation term on the
  strain-rate coupling, and the ordering of the two lags. Does not cover the existence theory of the
  initial-boundary-value problem. \\
Green--Naghdi type~II, in \cref{sec:hyperbolic} & Green and Naghdi (1993), Eq.~(3.16) integrated once &
  Covers the source-free form with a zero integration constant. Does not cover type~III, nor the entropy
  flux component relation, which is left open in the register. \\
Three-phase-lag law, \eqref{eq:tpl} & Roy Choudhuri (2007), Eqs.~(3)--(4) &
  Covers the first-order three-lag statement with its thermal-displacement-gradient term. Does not cover
  second-order truncations. \\
DPL stability ordering $\tau_T>\tau_q/2$ & Quintanilla (2003) &
  Covers a spectral (eigenvalue) stability statement for the linearized problem. Does \emph{not} rest on
  an energy functional, and says nothing about the nonlinear or finite-amplitude case. \\
Thermoelastic damping of the reed & Zener (1937) &
  Covers the attribution of the damping mechanism to Zener's reed calculation. Does not cover the
  plate and flexural extensions derived in \cref{app:disp}, which are re-derived here. \\
Second stiffness value in the piezothermoelastic bar example & The material data as stated with the
  example & The value is \emph{not} quoted from the earlier piezoelectric literature; the example states
  its own data, and the register marks the constitutive relation \eqref{eq:piezo} UNVERIFIED. \\
\bottomrule
\end{tabularx}
\end{table}

The last two rows show a deliberate pattern. Where a number could be borrowed from a source whose
equation was not checked, the number is either generated from stated data or omitted. A borrowed
number with an unchecked equation behind it is the most common way a review propagates an error it did not
commit.

\begin{summarybox}[Summary]
For linear homogeneous media, LS, GL and GN are covered by the one dispersion equation \eqref{eq:unified}; they differ in which coupling term carries a time-derivative correction. LS has finite thermal speed $c_T$ from flux relaxation, and GL has it from thermal inertia $\tau_0\ddot\theta$. GN-II is non-dispersive and energy-conserving. GN-III is wave-like at low frequency, with speed $c_*$, but diffusive at high frequency, so it has no sharp front and its phase speed is unbounded there. Front speeds obey \eqref{eq:fronts} and repel when $c_T\approx c_1$.
\end{summarybox}
